\documentclass[11pt, a4paper, oneside]{article}
\usepackage[ngerman]{babel}
\usepackage{worksheet}

\begin{document}
	\author{L. Bung}
	\title{Wiederholung \hspace{10cm} Polynomfunktionen}
	\subject{Mathematik}
	\class{FHR2}
	\maketitle
	
	\task{Polynomfunktionen zeichnen}
	
	Zeichnen Sie die Graphen der Funktionen.
	
	\begin{enumerate}[label=\alph*)]
		\item $f(x) = \frac{1}{20} x^4 - \frac{1}{4}x^3 - x$
		\item $g(x) = -x^4 + 4x^3 - 3x^2 + 2$
		\item $h(x) = -0.02x^4 + 0.5x^2 + x - 2$
	\end{enumerate}
	
	
	\task{Nullstellen berechnen}
	
	Berechnen Sie die Nullstellen der Funktionen.
	
	\begin{enumerate}[label=\alph*)]
		\item $f(x) = \frac{1}{20}x(x-2)^3(x+5)^2$
		\item $g(x) = x^4 - 5x^2 + 4$
		\item $h(x) = x^3 + 2x^2 + x$
	\end{enumerate}
	
	Geben Sie zu Teilaufgabe a) zusätzlich die Vielfachheit der Nullstellen an.
	
	
	\task{Globaler Verlauf von Polynomfunktionen}
	
	Beschreiben Sie den globalen Verlauf der Polynomfunktionen.
	
	\begin{enumerate}[label=\alph*)]
		\item $f(x) = 17x^5 - 3x^3 + x^2 - 15x + 27$
		\item $g(x) = 2x^4 - x^3 - 3x^2 - 4x - 15$
		\item $h(x) = -x^6 + 23x^4 - 2x^3 + 5x^2 + x + 18$
	\end{enumerate}
	
	
	\task{Symmetrien}
	
	Entscheiden und begründen Sie, ob die Funktionen
	\begin{itemize}
		\item[] $f(x) = 2x^4 + 4x^2 + 1$,
		\item[] $g(x) = x^3 + x^2 + x$ und
		\item[] $h(x) = -3x^5 + x^3 - x - 2$
	\end{itemize}
	Symmetrien haben.
	Falls ja, geben Sie zusätzlich die Art der Symmetrie an.
	
	
	\checkered[24cm]
\end{document}